
Requirements of Summary Section (500 words maximum)

In plain English, provide a summary we can use to identify the most suitable experts to assess your application.

We usually make this summary publicly available on external-facing websites, therefore do not include any confidential or sensitive information. Make it suitable for a variety of readers, for example:
opinion-formers
policymakers
the public
the wider research community

Guidance for writing a summary
Clearly describe your proposed work in terms of:
context
the challenge the project addresses
aims and objectives
potential applications and benefits

----------------

Title:
Addressing Global and Clinical Antimicrobial Resistance Using Genomics and Artificial Intelligence

Summary

Context 
The global epidemic of antimicrobial resistance (AMR) is one of the most urgent threats to modern medicine, affecting all countries and income levels. The World Health Organisation estimates that AMR was directly responsible for 1.27 million deaths globally and significantly contributed to 4.95 million deaths in 2019. The consequences extend far beyond clinical care: the World Bank estimates that AMR could result in US $1 - 3.4 trillion of loss to gross domestic product in 2030 and directly increase healthcare costs alone by US $1 trillion by 2050. This crisis is exacerbated by the continuous evolution of bacteria within patients, which acquire resistance genes, develop tolerance to antibiotics, and adapt to evade the immune system.

The Challenge
Despite the escalating threat of AMR, clinical tools for antibiotic prescribing have barely advanced since the early twentieth century. Standard laboratory tests report bacteria as simply "sensitive" or "resistant." This binary result is insufficient: it cannot predict if an apparently effective antibiotic will benefit a specific patient, nor explain why serious infections—both acute and chronic—often persist. This leads to treatment failure, recurring infections, and the wider spread of drug-resistant bacteria.

Aims and Objectives
This Professional Fellowship will harness cutting-edge bacterial genomics and artificial intelligence to build practical, clinically testable predictive tools that move beyond binary resistance testing. Focusing on Klebsiella pneumoniae, a critical, drug-resistant pathogen known for rampant gene swapping, this project builds on our genome-scale AI foundation model, Bacformer which reads an infections entire genetic code, through three connected objectives:
Understand Evasion: Conduct a global genomic study of K. pneumoniae to uncover the specific, hidden resistance and virulence mechanisms—missed by current tests, such as rampant gene swapping.
Develop AI Predictor: Use this new mechanistic knowledge to train and enhance Bacformer to predict resistance, tolerance, and virulence potential and understand gene swapping and insertions as K. pneumoniae shares traits with other species 
Creating a Clinical Tool: Build a practical hybrid clinical outcome model that combines the AI's genomic prediction score with a patient's medical history and clinical data from the SputaGen prospective study to accurately predict if an antibiotic course will fail, and if the patient's own bacteria will develop new resistance during treatment.
Potential Applications and Benefits
This project establishes a robust pipeline for developing and testing a new generation of genomics-informed diagnostics, guiding further refinement and clinical trials. The genomic "fingerprints" for each patient's pathogen will guide personalised antibiotic prescribing, helping to select antibiotics that will work, withhold those that will not, and predict the development of resistance. 

This directly addresses the global AMR crisis: at the individual level, it maximises the chance of successful treatment and extends antibiotic benefit and prevents the development of resistance within the patient's own microbiome; at the population level, it curtails the unnecessary use of ineffective drugs, reducing the spread of drug-resistant bacteria. 